\documentclass[12pt]{article}
\usepackage[margin=1in]{geometry}
\usepackage{amsmath,amssymb}
\usepackage{hyperref}
\usepackage{natbib}
\usepackage{graphicx}
\usepackage{float}

% Metadata
\title{The Rendering Universe: Causal Refinement as a Unified Framework for Cosmology}
\author{Patricio Fernando Bustos Cabrera\\
  \texttt{patriciob85@gmail.com}\\
  Maya Educaci\'on, Ecuador}
\date{July 2026}

\begin{document}
\maketitle

\begin{abstract}
We propose that cosmic expansion is not spatial stretching but {\bf causal refinement}: the dynamical increase of resolution $\mathcal{R}(t)$, measured in Planck pixels per side of a fixed Minkowski canvas of size 1. We present this as a {\bf phenomenological effective framework} whose macro-dynamics are consistent with several independent datasets. We show that (1) $\rho_{\rm vac} = M_{\rm Pl}^4 / \mathcal{R}^4$ resolves the 120-order vacuum catastrophe up to a factor $\sim40$ (as an IR/holographic regularization of gravitating vacuum energy), (2) a direct $\mathcal{R}\propto t^{\beta}$ likelihood against the official DESI DR2 BAO data (13 points, full covariance) gives $\beta_0 = 0.055^{+0.045}_{-0.050}$, resolving an earlier estimator tension, though BAO alone prefer this over $\Lambda$CDM by only $\Delta\chi^2\approx1$, (3) a pilot SPARC comparison shows a two-parameter cored profile fits better than NFW in 7/8 galaxies (though cored profiles outperform NFW independently of theory and the profile is not yet derived from the link-counting Postulate 6), (4) the variable-$\mathcal{R}$ account of the stacked ISW anomaly made a sharp, pre-registered prediction -- $\Delta T \approx -19\ \mu$K for the stack of the full DESIVAST DR1 void catalog against Planck -- which the measurement {\bf refuted}: $\Delta T = +1.2 \pm 1.5\ \mu$K, consistent with $\Lambda$CDM (repository experiment 17; falsifier~1 of Section~\ref{sec:falsifiers} fired, and the ISW sector is retracted as calibrated), and (5) JWST high-$z$ stellar mass densities are compatible with resolution sub-sampling only at $z \gtrsim 10$. A first-principles derivation of $\mathcal{R}(t,\mathbf{x})$ from causal-set dynamics remains the outstanding step; the surviving quantitative content is the BAO sector (2) and the numerically validated link-counting gravity of Postulate~6.
\end{abstract}

\section{Introduction}
\label{sec:intro}

The standard $\Lambda$CDM model successfully describes a wide range of cosmological observations but relies on several unexplained components: a cosmological constant fine-tuned to $10^{-120}$ in Planck units \citep{weinberg1989}, cold dark matter particles undetected after decades of searches \citep{bertone2018}, and an inflationary epoch driven by an unknown field \citep{guth1981}. The recent DESI DR2 results \citep{desi2025} showing $3.1$--$4.2\sigma$ preference for evolving dark energy ($w_0 > -1$, $w_a < 0$) add further tension.

We propose a simpler ontological shift: {\bf the universe is not expanding -- it is being rendered at increasing resolution.} This single hypothesis, inspired by the causal set approach to quantum gravity \citep{bombelli1987,sorkin2003}, replaces the three ad-hoc components with one mechanism.

\section{Theoretical Framework}
\label{sec:theory}

\subsection{The Central Postulate}
\label{sec:postulate}

{\bf Postulate 1 (Fixed Canvas).} Spacetime, at the fundamental level, is Minkowski flat:
\begin{equation}
ds^2_{\rm fundamental} = -c^2 dt^2 + dr^2 + r^2 d\Omega^2
\end{equation}
The universe is a manifold of total ``size'' 1 in suitable conformal units. It does not expand or contract.

{\bf Postulate 2 (Finite Resolution).} Observations are made with finite causal resolution $\mathcal{R}(t)$, the number of Planck pixels per side:
\begin{equation}
\mathcal{R}(t) \in [1, \infty)
\end{equation}
The resolution increases monotonically from $\mathcal{R} = 1$ at the Big Bang to $\mathcal{R}_0 \sim 10^{30}$ today, corresponding to $N_0 \sim \mathcal{R}_0^4 \sim 10^{120}$ causal elements (matching the holographic bound; \citealp{bousso2002}).

{\bf Postulate 3 (Resolution Redshift).} The cosmological redshift is not a stretching of space but a change in resolution between emission and reception:
\begin{equation}
1 + z = \frac{\mathcal{R}(t_0)}{\mathcal{R}(t_e)}
\end{equation}
This follows directly from comparing the same physical signal sampled at two different resolutions.

\subsection{Derivation from Causal Set Dynamics}
\label{sec:csg}

The framework is grounded in causal set theory (CST; \citealp{bombelli1987}). In CST, spacetime is a partially ordered set of $N$ elements. The number of elements $N$ scales as the spacetime volume in Planck units:
\begin{equation}
N \sim \frac{V}{\ell_P^4} \sim \mathcal{R}^4
\end{equation}

The causal set grows via {\bf Classical Sequential Growth (CSG)} \citep{rideout2000}. In CSG, new elements are added one at a time with transition probabilities that respect causality and discrete general covariance.

{\bf Derivation of $\mathcal{R}(t)$ dynamics.} The growth rate of causal set elements in Minkowski space ($V \propto t^4$) is:
\begin{equation}
\frac{dN}{dt} = \frac{d}{dt}\left(\frac{V(t)}{\ell_P^4}\right) = \frac{4V(t)}{\ell_P^4\,t}
\end{equation}
Using $N \sim \mathcal{R}^4$:
\begin{equation}
4\mathcal{R}^3 \frac{d\mathcal{R}}{dt} = \frac{4\mathcal{R}^4}{t} \quad\Rightarrow\quad \frac{d\mathcal{R}}{dt} = \frac{\mathcal{R}}{t}
\end{equation}
In a universe with matter, growth is modulated by the rate of causal connections per element:
\begin{equation}
\frac{d\mathcal{R}}{dt} = \frac{c}{\ell_P} \left(\frac{\rho_m(t)}{\rho_c(t)}\right)^{1/2}
\end{equation}
During matter and radiation domination, $\rho_m/\rho_c \approx 1$, recovering $\mathcal{R} \propto t$. During dark energy domination, $\rho_m/\rho_c \to 0$, slowing the refinement. The growth exponent $\beta \equiv d\ln\mathcal{R}/d\ln t$ evolves from $\beta \approx 1$ (early) to $\beta_0 \approx 0.05$--$0.15$ today; the current spread in $\beta_0$ reflects a real tension between two estimators, discussed openly in Section~\ref{sec:de}.

\subsection{The Effective Metric and Emergent Gravity}
\label{sec:metric}

If the fundamental metric is Minkowski but observations are made at resolution $\mathcal{R}(x)$, the effective metric was originally taken to be a conformal rescaling:
\begin{equation}
g_{\mu\nu}^{\rm (eff)}(x) = \eta_{\mu\nu} \frac{\ell_P^2}{\mathcal{R}(x)^2}
\label{eq:conformal-metric}
\end{equation}

{\bf Postulate 3$'$ (Effective Metric, amended 2026-07-04).} The render maps resolution to the effective metric \emph{anisotropically}: clocks respond as $1/\mathcal{R}$ and rulers as $\mathcal{R}$. In the weak field, with $\psi \equiv \ln(\mathcal{R}/\bar{\mathcal{R}})$,
\begin{equation}
ds^2 = -e^{-2\psi}\,c^2dt^2 + e^{+2\psi}\,d\mathbf{x}^2,
\label{eq:disformal}
\end{equation}
which yields $\gamma_{\rm PPN}=1$ by construction: gravitational light deflection, lensing and the Shapiro delay are recovered at linear order, and the Cassini bound $\gamma-1=(2.1\pm2.3)\times10^{-5}$ is satisfied. The conformal form of Eq.~(\ref{eq:conformal-metric}) is retained \emph{only} for the homogeneous cosmological background, where the two coincide up to a redefinition of the scale factor; as a local metric it gives $\gamma_{\rm PPN}=-1$ -- null geodesics are conformally invariant, so light would not bend at all -- and is observationally excluded at the $\sim10^5$ level. The amendment must apply universally (photons, gravitons and matter couple to the same $g_{\mu\nu}$), which keeps it consistent with the GW170817 bound $|c_{\rm GW}/c-1|<10^{-15}$. The $\mathcal{O}(\psi^2)$ sector (PPN $\beta$, perihelion precession) and the derivation of the anisotropic response from causal-set counting remain open problems; a covariant completion is deferred to future work.

The curvature formulas below are derived for the conformal metric of Eq.~(\ref{eq:conformal-metric}) and therefore apply to the {\bf homogeneous cosmological background only}; the local weak-field sector is governed by Postulate~3$'$, Eq.~(\ref{eq:disformal}). The Ricci curvature tensor for the conformally flat background metric is:
\begin{align}
R_{\mu\nu} &= \frac{2}{\mathcal{R}}\partial_\mu\partial_\nu\mathcal{R} - \frac{1}{\mathcal{R}}\Box\mathcal{R}\,\eta_{\mu\nu} \notag\\
&\quad - \frac{2}{\mathcal{R}^2}\partial_\mu\mathcal{R}\partial_\nu\mathcal{R} + \frac{1}{\mathcal{R}^2}(\partial\mathcal{R})^2\eta_{\mu\nu}
\end{align}
The Ricci scalar is:
\begin{equation}
R = -\frac{6}{\mathcal{R}}\Box\mathcal{R} + \frac{6}{\mathcal{R}^2}(\partial\mathcal{R})^2
\end{equation}

{\bf Limit of constant $\mathcal{R}$.} When $\partial_\mu\mathcal{R} = 0$, all curvature vanishes. The metric reduces to Minkowski locally.

{\bf Limit of slowly varying $\mathcal{R}$.} The Einstein tensor becomes:
\begin{equation}
G_{\mu\nu} \approx \frac{2}{\mathcal{R}}\bigl(\partial_\mu\partial_\nu\mathcal{R} - \eta_{\mu\nu}\Box\mathcal{R}\bigr)
\end{equation}
Identifying the effective stress-energy tensor of the resolution field:
\begin{equation}
T_{\mu\nu}^{(\mathcal{R})} \equiv \frac{1}{4\pi G\mathcal{R}}\bigl(\partial_\mu\partial_\nu\mathcal{R} - \eta_{\mu\nu}\Box\mathcal{R}\bigr)
\end{equation}
The cosmological constant emerges as $\Lambda = 3(\dot{\mathcal{R}}/\mathcal{R})^2_{\rm bg}$, matching the observed value.

\subsection{The Vacuum Catastrophe}
\label{sec:vacuum}

The quantum vacuum energy density, integrated to infinite momentum, diverges. We postulate that the vacuum energy \emph{that gravitates} is bounded by the coherence scale of the causal render:
\begin{equation}
\rho_{\rm vac} = \frac{1}{2\pi^2}\int_0^{2\pi/(\ell_P\mathcal{R})} k^2\sqrt{k^2+m^2}\,dk \approx \frac{2\pi^2 M_{\rm Pl}^4}{\mathcal{R}^4}
\label{eq:rhovac}
\end{equation}
with effective momentum bound $k_{\rm max} = 2\pi M_{\rm Pl}/\mathcal{R}$.\footnote{The factor $2\pi$ is absorbed into the definition of $\mathcal{R}$ in subsequent sections to keep notation simple.} Today, $\mathcal{R}_0 \sim 10^{30}$ yields $\rho_{\rm vac}^{\rm pred} \approx 10^{-120}\,M_{\rm Pl}^4$ (observed: $2.3\times10^{-122}\,M_{\rm Pl}^4$), removing the 120-order discrepancy up to a factor of $\sim 40$.

{\bf What this ansatz is and is not.} The bound $k_{\rm max} \sim M_{\rm Pl}/\mathcal{R}_0 \sim 10^{-11}$~eV must \emph{not} be read as a cutoff on laboratory quantum field theory, which is verified to TeV scales; modes above $k_{\rm max}$ exist and interact, but their zero-point energy is postulated not to source the effective $T_{\mu\nu}^{(\mathcal{R})}$. This is an infrared/holographic regularization of the \emph{gravitating} vacuum energy, in the spirit of the Cohen--Kaplan--Nelson bound $\rho_{\rm vac} \lesssim M_{\rm Pl}^2 L^{-2}$ for a region of size $L$ \citep{cohen1999} and of holographic entropy bounds \citep{bousso2002}: taking $L \sim R_H$ gives $\rho \sim M_{\rm Pl}^2 H_0^2 \sim 10^{-122}\,M_{\rm Pl}^4$, numerically consistent with our $\mathcal{R}^{-4}$ scaling. A first-principles derivation of which modes gravitate -- presumably from the causal-set structure, cf.\ the everpresent-$\Lambda$ program \citep{sorkin2005,ahmed2004} -- is an open problem of the framework; Eq.~(\ref{eq:rhovac}) is an ansatz whose independent support is the DESI-facing dynamics of Section~4, not a derivation from quantum field theory.

\subsection{Temperature as Sampling Noise}
\label{sec:temp}

The CMB temperature of $T_{\rm CMB} = 2.725$~K emerges naturally from the sampling noise of a discrete causal set. A first-principles derivation would require the full causal set dynamics; here we present only an empirical scaling ansatz. The observed temperature is reproduced by

\begin{equation}
T_{\rm CMB} \sim \frac{T_{\rm Planck}}{\mathcal{R}_0} \cdot \frac{1}{50}
\label{eq:cmb-ansatz}
\end{equation}

with $T_{\rm Planck} = \sqrt{\hbar c^5 / (G k_B^2)} \approx 1.4\times10^{32}$~K and $\mathcal{R}_0 \sim 10^{30}$, giving $T_{\rm Planck}/\mathcal{R}_0 \approx 140$~K, reduced by an empirical factor $\sim 1/50$ to match 2.725~K.\footnote{This is a purely empirical calibration; the factor $1/50$ is not derived from first principles. A complete treatment would compute the mode-counting on the 3-dimensional CMB surface, for which the 4-dimensional counting $T \sim T_{\rm Planck}/\mathcal{R}_0^2 \sim 10^{-28}$~K is far too small, suggesting that the relevant degrees of freedom are those of the surface rather than the bulk. This remains open for future work.} Deriving the exact prefactor from first principles is deferred to a dedicated treatment.

\section{Dark Energy as Quantization Noise}
\label{sec:de}

\subsection{Equation of State}
Since $\rho_{\rm DE} \propto \mathcal{R}^{-4}$, the equation of state is:
\begin{equation}
w(a) = -1 + \frac{4}{3}\frac{d\ln\mathcal{R}}{d\ln a}
\end{equation}
This follows from $w = -1 + (1/3)d\ln\rho/d\ln a$ and $\rho_{\rm DE} \propto \mathcal{R}^{-4}$. Converting the time-derivative of $\mathcal{R}$ to a scale-factor derivative requires an explicit $a(t)$ relation. In the matter-dominated era $a \propto t^{2/3}$, giving $d\ln\mathcal{R}/d\ln a = (3/2)\,d\ln\mathcal{R}/d\ln t \equiv (3/2)\,\beta$. With $\mathcal{R} \propto t^{\beta}$:
\begin{equation}
w_0 = -1 + 2\beta_0
\end{equation}
An earlier version of this section reported a factor $\sim$2--3 tension between two estimates of $\beta_0$: per-bin reconstruction of $\mathcal{R}(z)$ from the BAO distances gave $\beta_{\rm eff} \approx 0.02$--$0.07$, while mapping the published DESI $w_0w_a$CDM fit ($w_0=-0.727\pm0.067$; \citealp{desi2025}) through $w_0=-1+2\beta_0$ suggested $\beta_0\approx0.14$. We have now performed the direct likelihood fit that we had identified as the arbiter (repository experiment 15): the self-consistent model $\mathcal{R}\propto t^{\beta}$, $\rho_{\rm DE}\propto\mathcal{R}^{-4}$, fit to the official DESI DR2 release of 13 BAO measurements ($D_V, D_M, D_H$ over $r_d$) with the full published covariance, profiling over $\Omega_m$ and $H_0 r_d$. The result is $\beta_0 = 0.055^{+0.045}_{-0.050}$ (68\%, BAO alone; $0.040^{+0.040}_{-0.040}$ with a Planck-like $\Omega_m$ prior), with excellent goodness of fit ($\chi^2/{\rm dof}=0.92$). The reconstruction band lies inside the 68\% interval, while $\beta_0=0.14$ is excluded at 95\% once the $\Omega_m$ prior is applied: the apparent tension was an artifact of translating a different model ($w_0w_a$CDM) fit to a different data combination (BAO+CMB+SNe) through the matter-era conversion. Two honest caveats follow. First, BAO alone prefer $\beta>0$ over $\Lambda$CDM ($\beta=0$, nested) by only $\Delta\chi^2 \approx 1$: the framework is \emph{consistent with}, but not yet \emph{preferred by}, these data. Second, the exact present-day conversion is $w_0 = -1 + (4/3)\,\beta_0/(H_0 t_0)$, giving $w_0^{\rm eff} \approx -0.92$ at the best fit; the combined BAO+SNe likelihood required by falsifier~4 (Section~\ref{sec:falsifiers}) remains the outstanding test.
{\bf Caveat:} The conversion $a \propto t^{2/3}$ is standard FRW. In the strict Render framework, $1+z = \mathcal{R}_0/\mathcal{R}(t)$ replaces $a$ as the expansion parameter, and the relation $\mathcal{R} \propto t^{\beta}$ is the fundamental one. The $w$ parameterization is used here only to connect to DESI observations; a more fundamental prediction is $\mathcal{R}(t)$ directly.

\subsection{DESI DR2 Comparison}
DESI DR2 \citep{desi2025} measured BAO from 14 million galaxies ($0.1<z<4.2$). The combination DESI+CMB+SNe prefers evolving dark energy at $2.8$--$4.2\sigma$, with $w_0 > -1$ and $w_a < 0$. In the present framework this is interpreted phenomenologically: the strict relation $1+z=\mathcal{R}_0/\mathcal{R}$ gives the leading sampling-redshift map, while the effective DESI equation of state constrains the time-dependence of the refinement rate $\beta=d\ln\mathcal{R}/d\ln t$. A full likelihood analysis should fit $\mathcal{R}(t)$ directly rather than infer it through a standard $w_0$--$w_a$ parameterization.

\section{Dark Matter as Causal Aliasing}
\label{sec:dm}

{\bf Postulate 5 (Causal Aliasing).} What is called ``dark matter'' is a causal aliasing effect: the residual correlation of the primordial causal pixel that seeded a galaxy.

\subsection{Velocity Profile}
The causal halo's velocity profile:
\begin{equation}
V_{\rm halo}^2(R) = V_{\rm flat}^2 \frac{(R/R_c)^2}{1 + (R/R_c)^2}
\label{eq:causal-halo}
\end{equation}
For $R \ll R_c$: $V_{\rm halo} \propto R$; for $R \gg R_c$: $V_{\rm halo} \to V_{\rm flat}$.

\subsection{SPARC Test}
As a pilot test against the SPARC database \citep{lelli2016}, we fit eight representative galaxies with published rotation-curve decompositions:

\begin{table}[H]
\centering
\caption{Causal halo vs NFW fit for SPARC galaxies. Full 8-galaxy test sample.}
\begin{tabular}{lcccccc}
\hline
Galaxy & $N_{\rm pts}$ & $V_{\rm flat}$ & $R_c$ & $\chi^2_{\rm causal}$ & $\chi^2_{\rm NFW}$ & Winner \\
\hline
DDO154 & 12 & 50 & 3.2 & 1.10 & 10.10 & Causal \\
NGC2403 & 45 & 150 & 12.6 & 3.34 & 16.88 & Causal \\
NGC2841 & 49 & 235 & 9.9 & 0.67 & 2.15 & Causal \\
NGC3198 & 33 & 140 & 40.0 & 25.36 & 41.77 & Causal \\
NGC7331 & 14 & 125 & 40.0 & 19.46 & 23.78 & Causal \\
UGC09133 & 35 & 205 & 40.0 & 4.77 & 8.59 & Causal \\
NGC5907 & 19 & 175 & 40.0 & 12.80 & 19.78 & Causal \\
UGC02885 & 16 & 250 & 40.0 & 5.52 & 3.92 & NFW \\
\hline
\end{tabular}
\end{table}

The causal profile wins in 7 of 8 test galaxies. The single NFW winner (UGC02885) has a large $R_c$ indicating the causal halo has not yet transitioned to flat within the data range. This is a small test sample -- a full SPARC analysis with profile decomposition and Bayesian model comparison is deferred to future work.

\section{Cosmic Structure from Resolution}
\label{sec:structure}

{\bf Postulate 6 (Inhomogeneous Resolution via Link Counting).} The spatial variation of the resolution field is sourced by matter through the causal \emph{link} structure of the underlying causal set. Writing $\psi(x) \equiv \ln[\mathcal{R}(x)/\bar{\mathcal{R}}(t)]$, we define
\begin{equation}
\psi(x) = \mu \, K \sum_{s \,\in\, \mathrm{matter}} L(s, x),
\qquad K = \frac{1}{2\pi}\sqrt{\frac{\rho_c}{6}},
\label{eq:post6-links}
\end{equation}
where $L(s,x)=1$ if $(s,x)$ is a link (a causal relation with no intermediate element), $K$ is the normalization of the discrete retarded propagator of Johnston \citep{johnston2008}, $\rho_c$ is the sprinkling density, and $\mu$ a coupling constant. This definition is intrinsically non-local -- links concentrate on the past light cone \citep{myrheim1978,bombelli1987} -- and inherits the program of recovering geometry from causal order and counting \citep{benincasa2010}.

Two properties make this the correct formulation. First, numerical sprinkling experiments (repository experiments 12--13) show that the link-counted $\psi$ reproduces the Newtonian kernel quantitatively: potential $\propto 1/r$ (measured exponent $-0.984\pm0.031$), force $\propto 1/r^2$ ($-2.03\pm0.23$), static limit, and Johnston normalization verified to 1--3\% with no fitted parameters. Second, the same experiments show that the naive alternative -- counting all elements of the causal past interior -- fails: it has no static limit and yields a distance-independent force ($1/r^2$ excluded at $>500\sigma$), and is therefore \emph{discarded} as a definition of $\mathcal{R}$. Since $\psi > 0$ near mass, $\mathcal{R}$ is locally enhanced and proper time $d\tau = (\ell_P/\mathcal{R})\,dt$ runs slow near mass, with the correct sign for gravitational redshift. An earlier local form of this postulate, $\partial\mathcal{R}/\partial t \propto \rho(\mathbf{x},t)/\rho_c(t)$, is retained only as the homogeneous background growth law; as a local sourcing rule it cannot reproduce potential-dependent effects (e.g.\ tower gravitational redshift experiments, where local density is nearly equal at both heights) and is superseded by Eq.~(\ref{eq:post6-links}).

\subsection{JWST: Massive Galaxies at High Redshift}
The stellar mass density \citep{labbe2023,harvey2024,boylankolchin2023} follows from the resolution scaling. From Postulate 3 ($1+z = \mathcal{R}_0/\mathcal{R}$):
\begin{equation}
{\rm SMD}(z) \approx {\rm SMD}(0)
\left(\frac{\mathcal{R}(z)}{\mathcal{R}_0}\right)^3
= {\rm SMD}(0)\,(1+z)^{-3}
\label{eq:smd}
\end{equation}

\begin{table}[H]
\centering
\caption{JWST stellar mass density vs redshift.}
\begin{tabular}{cccc}
\hline
$z$ & JWST (log) & $\Lambda$CDM & This work \\
\hline
7.0 & $10^{7.36}$ & $10^{7.15}$ & $10^{6.27}$ \\
10.5 & $10^{6.20}$ & $10^{5.50}$ & $10^{5.91}$ \\
12.5 & $10^{5.68}$ & $10^{4.80}$ & $10^{5.76}$ \\
\hline
\end{tabular}
\end{table}

We note the crossover openly: the $(1+z)^{-3}$ sub-sampling scaling outperforms the $\Lambda$CDM expectation only at $z \gtrsim 10$ (0.3~dex vs 0.9~dex deviation at $z=12.5$), while at $z \approx 7$ it \emph{under}predicts the observed SMD by an order of magnitude, worse than $\Lambda$CDM. The natural reading is that Eq.~(\ref{eq:smd}) is at best the high-redshift limiting behavior, with standard hierarchical growth dominating by $z \sim 7$; a model combining both regimes is required before this can be claimed as evidence, and we list it as such rather than as a confirmed success.

\subsection{The Integrated Sachs-Wolfe Puzzle}
The ISW arises from variable $\mathcal{R}$ along the photon path \citep{sachs1967}. The cross-correlation ISW signal is also consistent with this picture \citep{krolewski2022}:
\begin{equation}
\frac{\Delta T}{T_{\rm CMB}} = \int_{\rm LOS} \frac{\dot{\mathcal{R}}}{\mathcal{R}} \frac{\delta\mathcal{R}(l)}{\bar{\mathcal{R}}} \frac{dl}{c}
\end{equation}

{\bf Empirical calibration.} Using the Hansen+ \cite{hansen2025} measurement at $z<0.03$:
\begin{equation}
\frac{\Delta T}{T_{\rm CMB}} = A\,\frac{L}{R_H}\,e^{-z/z_0},\quad A = 0.0024,\ z_0 = 0.30
\end{equation}

\begin{table}[H]
\centering
\caption{ISW stacking measurements vs prediction.}
\begin{tabular}{lcccc}
\hline
Experiment & $z$ & Observed ($\mu$K) & Predicted ($\mu$K) & $\Lambda$CDM \\
\hline
Granett+ 2008 \citep{granett2008} & 0.50 & $9.6\pm2.2$ & 7.0 & $\sim2$ \\
Kovacs+ 2021 \citep{kovacs2021} & 0.55 & $8.0\pm3.0$ & 5.9 & $\sim2$ \\
Hansen+ 2025 \citep{hansen2025} & 0.01 & $30\pm9$ & 28.7 & $\sim2$ \\
\hline
\end{tabular}
\end{table}

\section{Black Holes as Local $\mathcal{R}=1$}
\label{sec:bh}

{\bf Postulate 7 (Black Hole Pixel).} Inside a black hole horizon, the local resolution collapses to $\mathcal{R}=1$ -- a single causal pixel. The entropy $S = A/(4\ell_P^2)$ counts the elements on the horizon. Information is not lost but maximally compressed; Hawking radiation decompresses it over the evaporation time.

\section{The DESI $\times$ Planck Void Stacking Test: Prediction Refuted}
\label{sec:prediction}

The original version of this section predicted $\Delta T \approx 14\ \mu$K (in magnitude) for void stacking at $z\sim0.3$, $\sim7\times$ larger than $\Lambda$CDM's $\sim2\ \mu$K. {\bf We performed this test} (repository experiment 17): the full DESIVAST DR1 VoidFinder catalog (1{,}489 non-edge voids, $0.055<z<0.234$, median $R_{\rm eff}=22$~Mpc) was stacked against the Planck SMICA map with a compensated top-hat filter at each void's rescaled angular radius. The analysis was pre-registered and frozen (with author and independent audit) before unblinding; the null distribution (1{,}000 random-position sets) and a $-10\ \mu$K injection test validated the pipeline before any real void position was examined.

{\bf Result:}
\begin{equation}
\Delta T_{\rm stack} = +1.24 \pm 1.46\ \mu\text{K} \qquad (460\ \text{voids with valid patches}),
\end{equation}
consistent with zero and with $\Lambda$CDM. The frozen framework predictions -- $-18.5\ \mu$K (empirical calibration) and $-61.7\ \mu$K (cumulative-deficit mechanism of experiment 16) -- are excluded at $\sim$13$\sigma$ and $\sim$41$\sigma$ respectively. Robustness splits (NGC only; $R_{\rm eff}$ above the median) are also null. Falsifier~1 of Section~\ref{sec:falsifiers} has therefore {\bf fired}: the variable-$\mathcal{R}$ account of the stacked ISW anomaly is refuted as calibrated. Because the coupling $A$ was calibrated against extreme-supervoid measurements \citep{granett2008,hansen2025}, this null result on an unselected catalog also supports the interpretation that those anomalies reflect a-posteriori selection effects rather than a universal coupling; consequently the $\beta_{\rm ISW}$--$\beta_{\rm BAO}$ consistency reported in experiment 16 loses its evidential basis. Any revival of the ISW sector must be accompanied by a new pre-registered prediction (e.g.\ for an independent supervoid-selected sample) rather than a reinterpretation of this outcome.

\section{Discussion}
\label{sec:discussion}

\subsection{Relation to Existing Alternatives}
Causal set theory provides the closest technical precedent for this proposal: spacetime discreteness is encoded by a partial order, the number of elements estimates spacetime volume, and classical sequential growth gives a natural language for causal refinement \citep{bombelli1987,rideout2000,dowker2013}. The ``everpresent $\Lambda$'' program similarly argues that discreteness-induced fluctuations of causal-set volume can keep an effective cosmological term near the Hubble scale without a conventional $10^{-120}$ fine-tuning \citep{sorkin2005,ahmed2004}. The Render framework should be read as a coarse-grained phenomenological extension of that idea: $\mathcal{R}(t,\mathbf{x})$ is an effective macroscopic variable whose dynamics remain to be derived from an action or stochastic growth law.

The initial condition $\mathcal{R}=1$ has a precise counterpart in cyclic causal-set cosmology: a \emph{post} -- a single element through which every causal path passes -- separates successive cosmic cycles in CSG models, and Sorkin's ``cosmic renormalization'' shows that the effective dynamical parameters are renormalized at each passage through a post, providing a mechanism for parameter naturalness without fine-tuning \citep{sorkin2000,dowkerzalel2017}; observables for such cyclic models have recently been constructed \citep{dowkerzalel2023}. A speculative extension of the present framework would identify the $\mathcal{R}\to\infty$ asymptotics of one cycle with the $\mathcal{R}=1$ post of the next, in structural analogy with the conformal crossover of Penrose's Conformal Cyclic Cosmology \citep{gurzadyanpenrose2010}. We note, however, that CCC's proposed CMB signatures (low-variance concentric circles) were not confirmed by independent reanalyses \citep{moss2011}, and we do \emph{not} adopt the cyclic identification here: the framework as formulated requires only a single monotonic epoch, and none of its falsifiers (Section~\ref{sec:falsifiers}) depends on it.

\begin{table}[H]
\centering
\caption{Comparison with existing theories.}
\begin{tabular}{lll}
\hline
Theory & Relationship & Key difference \\
\hline
Causal sets \citep{bombelli1987,dowker2013} & Foundation & $\mathcal{R}(t)$ as coarse-grained variable \\
Everpresent $\Lambda$ \citep{sorkin2005,ahmed2004} & $\Lambda\!\propto\!1/\sqrt{V}$ & $\mathcal{R}^4$ scaling mechanism \\
Emergent gravity \citep{verlinde2017} & $a_0 = cH_0$ & Sampling, not entropy \\
MOND \citep{milgrom1983} & $a_0 = cH_0/2\pi$ & $a_0$ as render refresh rate \\
\hline
\end{tabular}
\end{table}

\subsection{Limitations}
\begin{enumerate}
\item The dynamics of $\mathcal{R}(t)$ require a first-principles derivation from the causal set action or a stochastic growth law.
\item The transition from $\mathcal{R}$ inhomogeneity to Einstein's equations requires a full action principle.
\item {\bf The ISW coupling derivation is void following the experiment-17 refutation.} The analysis below is retained for the record, but its premise -- that $A=0.0024$ measured a universal coupling -- did not survive the pre-registered stacking test (Section~\ref{sec:prediction}), which found no anomalous signal in an unselected void catalog. The $\beta_{\rm ISW}$--$\beta_{\rm BAO}$ consistency therefore no longer constitutes evidence. Original text follows. Repository experiment 16 shows that modeling a void as a region of \emph{cumulative refinement deficit}, $\delta\mathcal{R}/\bar{\mathcal{R}} = \gamma\beta|\delta_0| I(z)$ with $I(z)=\int D\,dt/t$, yields $A = \gamma\beta^2|\delta_0| I(0)/(H_0t_0)$ with no new parameters; reproducing the calibrated $A=0.0024$ requires $\beta_{\rm ISW}=0.045$, inside the 68\% BAO interval of Section~\ref{sec:de} ($\beta_{\rm BAO}=0.055^{+0.045}_{-0.050}$) -- two independent datasets pointing to the same refinement rate. Three caveats: (i) the redshift dependence is \emph{not} reproduced (the mechanism predicts mild growth with $z$, while the data decay as $e^{-z/z_0}$; $z_0=0.30$ remains empirical, plausibly encoding void-catalog selection and $\delta_0(z)$ evolution); (ii) the mechanism is the cumulative cosmological deficit, \emph{not} the static link potential of Postulate~6, which underpredicts the effect by four orders of magnitude -- the theoretical consistency between the two sectors is an open requirement; (iii) the sourcing exponent $\gamma\in[1/2,1]$ is not yet fixed by the theory.
\item The SPARC comparison is currently an eight-galaxy pilot test; a publishable claim requires the full sample, measurement errors, nuisance parameters, and Bayesian or information-criterion model comparison. Moreover, the causal-halo profile of Eq.~(\ref{eq:causal-halo}) is a \emph{cored} profile, and cored profiles are known to outperform NFW on dwarf and low-surface-brightness rotation curves independently of any specific theory (the core--cusp problem; \citealp{deblok2010}); the 7/8 result therefore does not by itself discriminate this framework from other cored alternatives.
\item {\bf The causal-halo profile is not derived from Postulate 6.} The link-counting definition of $\mathcal{R}$ sourced by baryons reproduces the Newtonian kernel (experiments 12--13), and Newtonian gravity from baryons alone does not produce flat rotation curves. Flat curves therefore require additional physics not yet present in the framework (e.g.\ inhomogeneous-sprinkling effects or back-reaction of $\mathcal{R}$ gradients), and Eq.~(\ref{eq:causal-halo}) must currently be read as an empirical fitting function pending derivation.
\item {\bf The Bullet Cluster.} Any resolution field sourced solely by baryonic matter faces the same objection as baryon-sourced modified gravity: in merging clusters the weak-lensing mass peaks separate from the dominant baryonic component (the X-ray gas) and track the galaxies \citep{clowe2006}. Explaining this requires a dynamical rule for how collisions affect link counting differently for gas and galaxies -- no such rule is currently formulated. Until it is, the framework has no account of merging-cluster lensing.
\item Quantum mechanics is not addressed.
\end{enumerate}

\subsection{What Would Falsify This Framework}
\label{sec:falsifiers}

A framework whose central variable parameterizes the whole cosmic history risks becoming unfalsifiable by reinterpretation. To guard against this, we state in advance the observations that would refute the framework \emph{as formulated in this paper}, without appeal to new parameters or post-hoc modifications:

\begin{enumerate}
\item {\bf DESI$\times$Planck void stacking at $z\sim0.3$ yields $\Delta T \approx 2\ \mu$K.} The flagship prediction of Section~\ref{sec:prediction} was $\approx14\ \mu$K in magnitude, a factor $\sim$7 above $\Lambda$CDM. A stacking measurement consistent with $\Lambda$CDM refutes the variable-$\mathcal{R}$ account of the ISW anomaly. {\bf STATUS (2026-07-02): FIRED.} The pre-registered measurement (Section~\ref{sec:prediction}; repository experiment 17) gave $+1.24\pm1.46\ \mu$K, consistent with $\Lambda$CDM; the framework predictions were excluded at $\gtrsim$13$\sigma$. The ISW sector is retracted as calibrated. We record the outcome here rather than deleting the falsifier, as the falsification discipline requires.
\item {\bf A confirmed phantom equation of state, $w_0 < -1$.} Postulate~2 requires $\mathcal{R}(t)$ to grow monotonically, hence $\beta_0 > 0$ and $w_0 = -1 + 2\beta_0 \geq -1$ with no adjustable freedom. A robust cosmological convergence on $w_0 < -1$ (or a confirmed phantom crossing) falsifies the dark-energy mechanism outright. This is the cleanest falsifier: it depends on no fitted parameter of the framework.
\item {\bf Direct detection of a dark matter particle} with the relic abundance required by $\Lambda$CDM. Postulate~5 attributes the missing-mass phenomenology to causal aliasing rather than particles; a confirmed particle detection removes its motivation and the framework's dark-matter sector with it.
\item {\bf An irreconcilable $\beta_0$.} Section~\ref{sec:de} reports a factor $\sim$2--3 tension between the BAO reconstruction of $\beta_0$ ($\approx0.02$--$0.07$) and the value implied by the DESI parametric fit ($\approx0.14$). We commit to the direct $\mathcal{R}(t)$ likelihood analysis as the arbiter: if no single monotonic $\mathcal{R}(t)$ fits the combined BAO+SNe data, the quantitative dark-energy sector fails even though the qualitative quadrant ($w_0>-1$, $w_a<0$) matches.
\item {\bf Energy-dependent photon dispersion.} The framework as formulated assumes the underlying sprinkling preserves local Lorentz invariance \citep{bombelli1987}, so photon propagation speed must be independent of energy and of local $\mathcal{R}$ to current experimental precision. Fermi-LAT gamma-ray-burst bounds already constrain linear Lorentz violation beyond the Planck scale; if any future refinement of the render dynamics predicts energy-dependent arrival times, those bounds apply immediately. Conversely, a confirmed detection of Lorentz-violating dispersion would require a fundamental revision, since the fixed Minkowski canvas (Postulate~1) provides no mechanism for it.
\end{enumerate}

Each rescue of the framework from one of these outcomes would have to generate new testable predictions of its own; rescues that merely reinterpret a failure are not admissible. We adopted this discipline internally as well: the link-counting formulation of Postulate~6 replaced an earlier local form precisely because pre-registered numerical experiments (repository experiments 12--13) refuted the original mechanism, and the negative results are published in the repository alongside the positive ones.

\section{Conclusion}
\label{sec:conclusion}

We have outlined a single hypothesis -- the universe undergoes causal refinement rather than spatial expansion -- and subjected it to its own falsifiability program. The outcome to date is mixed and is reported without adjustment. Refuted: the variable-$\mathcal{R}$ account of the stacked ISW anomaly, by our own pre-registered DESI$\times$Planck measurement (Section~\ref{sec:prediction}), which came out consistent with $\Lambda$CDM. Surviving: the direct $\mathcal{R}\propto t^{\beta}$ likelihood against DESI DR2 BAO ($\beta_0 = 0.055^{+0.045}_{-0.050}$, consistent with but not preferred over $\Lambda$CDM), and the numerically validated link-counting formulation of Postulate~6, which reproduces the Newtonian kernel from causal structure with no fitted parameters. Open: the first-principles derivation of $\mathcal{R}(t,\mathbf{x})$ from a causal-set action, the dark-matter sector, and the reconciliation of the sampling and refinement roles of $\mathcal{R}$. The framework is not a replacement for $\Lambda$CDM. We submit that its chief contribution, beyond the surviving sectors, is methodological: a speculative cosmology can and should ship with its own falsifiers, and record their outcomes when they fire.

\section*{Acknowledgments}
The author thanks the DESI Collaboration, the SPARC team, the Planck Collaboration, and Hansen et al. for making their data publicly available. Large language models (LLMs) were used as writing and mathematical formalization assistants; all scientific content is the author's own.

\section*{Data Availability}
All data and code are publicly available at \url{https://github.com/Editorenbici/rendering-universe}.

\begin{thebibliography}{30}

\bibitem[Weinberg(1989)]{weinberg1989} S.~Weinberg, Rev. Mod. Phys. {\bf 61}, 1 (1989).

\bibitem[Bertone and Hooper(2018)]{bertone2018} G.~Bertone and D.~Hooper, Rev. Mod. Phys. {\bf 90}, 045002 (2018).

\bibitem[Guth(1981)]{guth1981} A.~Guth, Phys. Rev. D {\bf 23}, 347 (1981).

\bibitem[DESI Collaboration(2025)]{desi2025} DESI Collaboration, Phys. Rev. D {\bf 112}, 083515 (2025).

\bibitem[Bombelli et al.(1987)]{bombelli1987} L.~Bombelli, J.~Lee, D.~Meyer and R.~Sorkin, Phys. Rev. Lett. {\bf 59}, 521 (1987).

\bibitem[Sorkin(2003)]{sorkin2003} R.~Sorkin, in {\it Lectures on Quantum Gravity}, pp.~305--327, Springer (2003).

\bibitem[Bousso(2002)]{bousso2002} R.~Bousso, Rev. Mod. Phys. {\bf 74}, 825 (2002).

\bibitem[Rideout and Sorkin(2000)]{rideout2000} D.~Rideout and R.~Sorkin, Phys. Rev. D {\bf 61}, 024002 (2000).

\bibitem[Sorkin(2005)]{sorkin2005} R.~Sorkin, in {\it Discrete Approaches to Quantum Gravity}, World Scientific (2005).

\bibitem[Ahmed et al.(2004)]{ahmed2004} M.~Ahmed, S.~Dodelson, P.~B.~Greene and R.~D.~Sorkin, Phys. Rev. D {\bf 69}, 103523 (2004).

\bibitem[Dowker(2013)]{dowker2013} F.~Dowker, Gen. Relativ. Gravit. {\bf 45}, 1651 (2013).

\bibitem[Lelli et al.(2016)]{lelli2016} F.~Lelli, S.~McGaugh and J.~Schombert, AJ {\bf 152}, 157 (2016).

\bibitem[Granett et al.(2008)]{granett2008} B.~Granett, M.~Neyrinck and I.~Szapudi, ApJ {\bf 683}, L99 (2008).

\bibitem[Kovacs et al.(2021)]{kovacs2021} A.~Kovacs et al., MNRAS {\bf 510}, 216 (2021).

\bibitem[Hansen et al.(2025)]{hansen2025} F.~Hansen et al., A\&A (accepted), arXiv:2506.08832 (2025).

\bibitem[Krolewski and Ferraro(2022)]{krolewski2022} A.~Krolewski and S.~Ferraro, JCAP {\bf 04}, 033 (2022).

\bibitem[Verlinde(2017)]{verlinde2017} E.~Verlinde, SciPost Phys. {\bf 2}, 016 (2017).

\bibitem[Milgrom(1983)]{milgrom1983} M.~Milgrom, ApJ {\bf 270}, 365 (1983).

\bibitem[Labb\'e et al.(2023)]{labbe2023} I.~Labb\'e et al., Nature {\bf 616}, 266 (2023).

\bibitem[Harvey et al.(2024)]{harvey2024} T.~Harvey et al., ApJ {\bf 978}, 89 (2024).

\bibitem[Boylan-Kolchin(2023)]{boylankolchin2023} M.~Boylan-Kolchin, Nature Astronomy {\bf 7}, 731 (2023).

\bibitem[Sachs and Wolfe(1967)]{sachs1967} R.~Sachs and A.~Wolfe, ApJ {\bf 147}, 73 (1967).

\bibitem[Johnston(2008)]{johnston2008} S.~Johnston, Class. Quantum Grav. {\bf 25}, 202001 (2008), arXiv:0806.3083.

\bibitem[Myrheim(1978)]{myrheim1978} J.~Myrheim, CERN preprint TH-2538 (1978).

\bibitem[Benincasa and Dowker(2010)]{benincasa2010} D.~M.~T.~Benincasa and F.~Dowker, Phys. Rev. Lett. {\bf 104}, 181301 (2010), arXiv:1001.2725.

\bibitem[Cohen et al.(1999)]{cohen1999} A.~G.~Cohen, D.~B.~Kaplan and A.~E.~Nelson, Phys. Rev. Lett. {\bf 82}, 4971 (1999).

\bibitem[de Blok(2010)]{deblok2010} W.~J.~G.~de Blok, Adv. Astron. {\bf 2010}, 789293 (2010).

\bibitem[Clowe et al.(2006)]{clowe2006} D.~Clowe et al., ApJ {\bf 648}, L109 (2006).

\bibitem[Sorkin(2000)]{sorkin2000} R.~D.~Sorkin, Int. J. Theor. Phys. {\bf 39}, 1731 (2000), gr-qc/0003043.

\bibitem[Dowker and Zalel(2017)]{dowkerzalel2017} F.~Dowker and S.~Zalel, C. R. Physique {\bf 18}, 246 (2017), arXiv:1703.07556.

\bibitem[Dowker and Zalel(2023)]{dowkerzalel2023} F.~Dowker and S.~Zalel, Class. Quantum Grav. {\bf 40}, 155015 (2023), arXiv:2212.01149.

\bibitem[Gurzadyan and Penrose(2010)]{gurzadyanpenrose2010} V.~G.~Gurzadyan and R.~Penrose, arXiv:1011.3706 (2010).

\bibitem[Moss et al.(2011)]{moss2011} A.~Moss, D.~Scott and J.~P.~Zibin, JCAP {\bf 04}, 033 (2011), arXiv:1012.1305.

\end{thebibliography}

\end{document}
